% !TeX root = ../main.tex
\chapter{Procedimentos de Teste e Validação do Projeto}
\label{cap:testes}

\section{Testes de caixa preta}

Os testes exercitam a CLI pela interface disponível ao desenvolvedor. Para os cenários de falha controlada, será usado um provedor substituto de LLM, com respostas previamente definidas. Os resultados desses testes verificam o protótipo. As métricas experimentais serão calculadas sobre execuções com o modelo escolhido para o estudo.

\begin{table}[H]
\centering
\caption{Cenários de caixa preta.}
\label{tab:cp}
\small
\begin{tabular}{@{}p{5.1cm}p{9.2cm}@{}}
\toprule
\textbf{Entrada ou cenário} & \textbf{Resultado esperado} \\
\midrule
SARIF válido ou manifesto com a mesma localização & Identificação do mesmo endpoint e registro da origem da entrada. \\
Alerta sem localização válida ou fora do recorte & Diagnóstico explícito, com motivo da interrupção. \\
Resposta com diff inválido ou arquivo protegido & Proposta recusada antes da compilação e checkout original intacto. \\
Proposta que preserva a operação e passa nos testes de autorização & Resultados de segurança e funcionalidade aprovados. Aderência registrada em separado. \\
Proposta que bloqueia um usuário autorizado ou mantém acesso indevido & Reprovação registrada no critério correspondente. \\
Contexto insuficiente ou falha no ambiente & Incerteza, abstenção ou erro de infraestrutura identificado no relatório, conforme o caso. \\
\bottomrule
\end{tabular}
\par\smallskip{\footnotesize Fonte elaborada pelos autores.}
\end{table}

\section{Testes de caixa branca}

A normalização será testada com regras mapeadas, regiões ambíguas e caminhos externos. A análise com Roslyn será testada com controladores, grupos de Minimal APIs, políticas padrão e de \textit{fallback}, metadados \texttt{AllowAnonymous} e verificações por recurso. Os testes verificarão a resolução dos símbolos e a sinalização de relações que permanecem indeterminadas.

O recuperador será testado com endpoints de operações semelhantes sujeitos a regras diferentes. Os testes verificarão a ordenação, as referências de origem e o limite de contexto. Os testes de geração verificarão o conteúdo enviado nas condições A, B e C, a ausência dos arquivos protegidos na entrada e o tratamento de respostas inválidas. O validador será testado com alterações fora do escopo, falha de compilação, regressões funcionais, falhas de autorização e interrupções durante a execução. Ao final de cada execução, os testes verificarão o descarte da cópia de trabalho e a integridade do checkout original.

\section{Critérios de avaliação}
\label{sec:criterios}

\begin{enumerate}
    \item \textbf{Segurança.} Os testes verificam a matriz esperada de acesso. Uma operação negada não deve produzir alteração de estado ou exposição de dados. Perfis e recursos de teste são definidos por caso.
    \item \textbf{Funcionalidade.} O código deve compilar, as regressões funcionais devem passar e as identidades autorizadas devem conservar o comportamento esperado, incluindo os efeitos da operação.
    \item \textbf{Aderência.} O \textit{diff} deve reutilizar a política, o requisito ou o serviço aplicável ao caso, no nível apropriado, com evidência no repositório. A revisão examinará regras inventadas, verificações duplicadas, alterações fora do escopo e divergência em relação à regra de negócio documentada.
\end{enumerate}

Os testes de integração do ASP.NET Core poderão usar \texttt{WebApplicationFactory} para iniciar a aplicação em um servidor de teste \cite{aspnettesting}. Antes do experimento, o avaliador verificará que o conjunto de testes distingue o comportamento inicial do comportamento esperado e que a correção de referência preserva a operação autorizada. Os testes e a referência serão mantidos fora do contexto de geração. A revisão de aderência usará uma rubrica definida antes das execuções e omitirá a condição experimental do avaliador, quando possível.

\section{Desenho experimental e métricas}

Cada caso elegível será executado nas três condições da Tabela~\ref{tab:condicoes}, sempre a partir do mesmo commit. Modelo, versão, parâmetros, instruções comuns, limite de saída e procedimento de validação serão fixados. Será solicitada uma proposta por condição, sem devolver ao modelo resultados dos testes. O orçamento de recuperação será calibrado no piloto e fixado antes de qualquer avaliação ampliada. Para B e C, o arquivo completo deve caber na entrada com a capacidade reservada para a recuperação em C. Casos que excederem esse limite serão registrados antes da comparação.

\begin{table}[H]
\centering
\caption{Condições de geração.}
\label{tab:condicoes}
\small
\begin{tabular}{@{}p{4.1cm}p{10.2cm}@{}}
\toprule
\textbf{Condição} & \textbf{Contexto de código enviado com o alerta} \\
\midrule
A --- Mínimo & Método ou manipulador do endpoint afetado. \\
B --- Arquivo & Arquivo completo que contém o endpoint. \\
C --- RAG & Mesmo arquivo completo de B e evidências de autorização recuperadas do repositório. \\
\bottomrule
\end{tabular}
\par\smallskip{\footnotesize Fonte elaborada pelos autores.}
\end{table}

A métrica principal será a fração de casos elegíveis em que a primeira proposta passa nos critérios de segurança, funcionalidade e aderência. A métrica de controle \textit{pass@1} será a fração que passa nos critérios de segurança e funcionalidade, independentemente da aderência. Abstenções, propostas inválidas e falhas do patch contam como insucessos. Uma falha de infraestrutura poderá motivar nova execução com a mesma configuração. Se persistir e impedir a comparação, o caso será informado separadamente como incompleto, com motivo e contagem. O denominador das taxas será o número de casos com as três condições avaliáveis.

Na condição C, será registrada a precisão dos trechos recuperados, medida pela proporção de trechos relevantes entre os retornados, e a cobertura das evidências essenciais anotadas antes da geração. Ausência de contexto será registrada separadamente. O relatório também informará tempo de recuperação, geração e validação, uso de tokens e custo quando disponível. A análise apresentará contagens, taxas e diferenças pareadas entre condições, com destaque para C em relação a B. Os resultados do piloto servirão para verificar o desenho, incluindo a existência de casos que de fato dependam de evidências fora do arquivo. As conclusões serão limitadas aos casos e ao modelo avaliados.
